% !TEX root = ../main.tex
\lecture{3}{2026-09-28}{Woche 3 Time Value of Money}

\section{Time Value of Money}

\subsection{Einführung}

\begin{definition}
    [Compound Interest oder Zinseszins] exponential funktion aus einer Anfangsmenge oder Present Value (PV) und wiederholter Zinsatz (k) auf neuer Wert der letzten Periode über eine Periodendauer (t) resultiert in Future Value (FV):
    \begin{equation}
        f(x) = G * Z^n
    \end{equation}
\end{definition}
\textbf{Fragen}
\begin{itemize}
    \item Wie wächst ein PV mit Zinsatz k in Periode t?
    \item Welcher PV muss heute investieren werden um mit k einen gewissen FV zu erreichen?
    \item Wie vie
\end{itemize}
\textbf{Gründe}
\begin{itemize}
    \item Geld kann angelegt/investiert werden
    \item Zukünftige Aufwertung ist unsicher
    \item Inflation mindert die Kaufkraft
\end{itemize}

\begin{definition}
    [\textit{diskrete} Unterjährige Verzinsung] wenn innerhalb einer Periode t eine wiederholte Verzinsung stattfindet.
    \begin{equation}
        FV = PV * (1+\frac{R}{m})^{t*m}
    \end{equation} 
    Wobei $R:$ unterjähriger Zinssatz oder Gesamtzinssatz, $m$: Anzahl Verzinsungen im Jahr.
\end{definition}

\begin{definition}
    [\textit{stetige} Verzinsung] form 
    \begin{equation}
        FV = K_0 \cdot e^{r_c \cdot T}
    \end{equation}
\end{definition}    

\begin{definition}
    [Kapitalwertmethode oder Net Present Value] 
    \begin{equation}
        NPV = - I_0 + \sum_{t=1}^{T}[\frac{CF_t}{(1 + k)^t}]
    \end{equation}
\end{definition}

Barwert oder NPV repräsentiert der Wert eines zinsfälligen Investitionsprojekts in einer \footnote{Ein Projekt mit NPV  = 0, rendiert nur so viel wert, wie die investition benötigt. Also keinen Profit.}